\documentclass{article}
\usepackage[T1]{fontenc}
\usepackage{lmodern}
\usepackage{graphicx}
\usepackage{amsmath,amssymb}
\usepackage[a4paper,margin=2cm]{geometry}
\usepackage[absolute,overlay]{textpos}
\setlength{\TPHorizModule}{1mm}
\setlength{\TPVertModule}{1mm}
\usepackage[indonesian]{babel}
\usepackage{hyperref}
\setlength{\emergencystretch}{2em}
% unit: Halaman 39

\begin{document}

% === Halaman 39 (llm) ===

\begin{itemize}
    \item Misalkan kita mengetahui pasangan potongan plaintext ($p$) dan cipherteks ($c$) yang berkoresponden.
    \item Lakukan serangan \textit{brute force} (atau \textit{exhaustive key search attack}) dengan mencoba semua kemungkinan kunci ($k$):
\end{itemize}

\vspace{5mm}

\begin{verbatim}
for (k=0; k<2^{56}; k++) {
    if (DES(k, p) == c) {
        printf("Key is %x\n", k);
        break;
    }
}
\end{verbatim}

\vspace{5mm}

\begin{itemize}
    \item Kunci diharapkan ditemukan sebelum $2^{55}$ iterasi.
\end{itemize}

\end{document}
